\section{Whittaker categories in families}\label{sec:whittaker}

The Whittaker coefficient functor will be used both in duality and in
families in which marked points collide. We construct its projection on
all D-modules, establish its compatibility with those families, and fix the
pairing and t-structure normalizations. The finite averaging calculation
below is the common source of these properties.

\subsection{The pole and Grassmannian models}

Let $H$ be one of the semisimple groups under consideration. Fix a Borel
$B=TN$ and the Cartan bundle $T_0$ inducing the bundle $P_0$ of
Section~\ref{sec:foundations}. Write
\[
 N_0=N^{T_0},\qquad \mathfrak n_0=\operatorname{Lie}(N_0),\qquad
 \delta_N=-\chi(X,\mathfrak n_0)=\dim\operatorname{Bun}_{N_0}.
\]
Here $\chi(X,-)$ denotes Euler characteristic. Each simple-root line of
$\mathfrak n_0$ is $\Omega_X$. The pinning and the sum of residues therefore
define a nondegenerate additive character, also denoted $\chi$, on the
unipotent loop groups. We use both $\chi$ and $-\chi$. Our gluing convention
sends the inside frame to the outside frame.

A parameter scheme $S$ carries a nonempty tuple $(x_i)_{i\in I}$ of points
of $X$; auxiliary parameters are allowed. The principal examples are the
powers $X^I$ and their coincidence strata. All operations in the parameter
scheme use the D-module conventions of Section~\ref{foundations:operations},
without a perverse shift unless one is displayed. The ordinary Ran space
is presented by the $X^I$, with transition maps the diagonals indexed by
surjections of nonempty finite sets.

Let $\mathcal Y\to S$ be the following pole stack. An object consists of an
$H$-bundle $P$ together with Pl\"ucker maps
\[
 T_0^\lambda\longrightarrow V_P^\lambda
 \qquad(\lambda\text{ a dominant weight of }H),
\]
regular away from the marked points, generically nonzero, and satisfying
the Cartan-product relations. All representations are those of the given
group $H$. Thus the definition uses the affine closure of $H/N$ and allows
defects; it imposes no quotient by Cartan scalings. Denote the forgetful map
by $p_{\mathcal Y}:\mathcal Y\to\operatorname{Bun}_H$ and give $\mathcal Y$
the pullback of the determinant twist $L_H^s$.

We write $\mathcal Y_m$ for closed pole bounds. Choose effective bounds for
a finite set of highest-weight generators, compatible with Cartan
multiplication; multiples of a coweight strictly positive in the coroot
span give a cofinal choice. At colliding markings these bounds add. Each
$\mathcal Y_m$ is locally QCA, and its map to
$\operatorname{Bun}_H\times S$ is schematic, separated, and of finite type:
after bounding poles, the finitely many maps are sections of vector bundles
satisfying equations and open conditions.

On each $\mathcal Y_m$ take the full subcategory of
$D_{p_{\mathcal Y}^*L_H^s}(\mathcal Y_m)$ defined by the strong Whittaker
condition: its objects are equivariant with character $\chi$ for
$N$-modifications away from the marked points and on the nondefective
locus. The determinant twist is split along these modifications by the
Borel filtration and the relative determinant identification. Restriction
between closed pole bounds respects this character descent. Define
$\mathcal W_s^+(S)$ to be the limit of these full subcategories under
$!$-restriction, inside the ambient category
\[
 D_{p_{\mathcal Y}^*L_H^s}(\mathcal Y)
 =\lim_m D_{p_{\mathcal Y}^*L_H^s}(\mathcal Y_m).
\]
The opposite category $\mathcal W_{-s}^-(S)$ is defined with $L_H^{-s}$
and $-\chi$. If the derived group is not simply connected, the usual
Drinfeld compactification imposes an additional closed defect condition.
Whittaker objects on the Pl\"ucker model used here are supported on that
closed substack \cite[Remark~4.1.2]{GaitsgoryWhit}; hence the two models
define the same Whittaker category. The projection and continuity properties, and the equivalent
colimit presentation, will be proved below. One may impose this condition by choosing an auxiliary point $z$, passing
to the open where $z$ is neither marked nor defective, framing the
reduction at $z$, and using the unipotent loops there. We verify below that
the resulting condition is independent of these choices.

The Grassmannian $\operatorname{Gr}_{P_0,S}$ parametrizes modifications of
$P_0$ at the marked points. Such a modification supplies the Pl\"ucker
maps and hence a morphism
\[
 \mu:\operatorname{Gr}_{P_0,S}\longrightarrow\mathcal Y.
\]
Its twisting line is the relative determinant: the pullback of $L_H$
divided by its fixed value at $P_0$. This line factorizes over disjoint
formal discs and has the unipotent splitting just specified.

\begin{proposition}[Comparison of Whittaker models]
\label{whittaker:model-comparison}
For a connected semisimple complex group $H$, a smooth complete curve $X$,
and the fixed half-canonical datum defining $T_0$, the functor $\mu^!$
identifies the pole-model Whittaker category with the strongly
$\chi$-equivariant D-module category on
$\operatorname{Gr}_{P_0,S}$ with its relative determinant twist. This holds
for every exponent $s$, with arbitrary D-modules in the marked-point
parameters, and is compatible with the multipoint and Ran presentations.
The assertion also holds with both signs reversed.
\end{proposition}

\begin{proof}
For the untwisted sheaf theory this is Gaitsgory's local/global comparison,
\cite[Theorems~5.5.2 and~7.2.3]{GaitsgoryWhit}. We explain why its proof
applies to our determinant twist, so that the distinction between a
character twist and a twisting line is explicit.

Let $\mathcal L=p_{\mathcal Y}^*L_H$ and form the frame torsor
$\mathcal L^\times\to\mathcal Y$. Pull every correspondence in the
comparison proof back from its target to this torsor. On an open
action correspondence, the unipotent and arc splittings transport this
frame to the source. On a compactified correspondence we pull back only
the target torsor; the open source-frame identification is not extended
as a trivialization over its boundary. This single Cartesian pullback
preserves the proper maps in the compactified action diagrams.

There are two ingredients in the comparison proof
\cite[Sections~6.1--6.6]{GaitsgoryWhit}: full faithfulness after Ran
uniformization, and the equivalence obtained by restricting to the unit
configuration in the Whittaker category. The geometric input to the first
is universal homological contractibility. This means precisely that the
corresponding pullback is fully faithful after every prestack base change;
see \cite[Appendix~A.1.8--A.1.11]{GaitsgorySemiInfiniteII}. It therefore
applies after the frame-torsor base change as well as with the marked-point
parameters retained.

The second ingredient uses unipotent averaging, arcs, and conjugates of
positive-depth subgroups. The determinant central extension is canonically
split on arcs, and the Borel filtration splits it on unipotent
modifications. These splittings restrict to the finite quotients used in
averaging and respect multiplication. For the subgroups conjugated by a
multiple of the half-sum of positive coroots, take the cofinal sequence of
even depths. The conjugating cocharacter then belongs to the actual
cocharacter lattice of $H$, since twice that half-sum is the sum of
coroots. Conjugation of the arc splitting is consequently defined: choose
a lift of the conjugating loop to the central extension, locally if
necessary; any two choices differ centrally and give the same result.

There is also a family compatibility to check at collisions. In the
notation of \cite[Section~6.6.4]{GaitsgoryWhit}, the group
$\mathfrak L_I^+(H)^j$ is the inverse image of the preceding one-point
subgroup under
$\mathfrak L_I^+(H)'\to\mathfrak L_x(H)$. On this group the relative
determinant extension is pulled back from the distinguished disc at $x$:
the primed group allows poles only on the distinguished graph
$E=S\times\{x\}$. More explicitly, let $t$ cut out $E$ and let $A_I$
be the ring completed along the entire marked divisor, including $E$. On a bounded
piece choose $m$ so that
$t^m\Lambda\subset g\Lambda\subset t^{-m}\Lambda$. The finite
lattice quotients defining the relative determinant are supported on a
thickening of $E$. For every $r$,
\[
 A_I/(t^r)\simeq\mathcal O_{S\times X}/(t^r)
             \simeq\mathcal O_S[[t]]/(t^r),
\]
because the ideal of the entire marked divisor is nilpotent modulo
$t^r$. Thus restriction to the distinguished disc preserves the quotient
complexes and their determinant identifications, including on collision
and nonreduced parameter schemes. The splitting therefore pulls back from the one-point
subgroup. It agrees on overlaps with the canonical arc splitting;
on pro-unipotent kernels uniqueness follows from the absence of
algebraic characters to $\mathbb G_m$. These constructions commute with
coordinate changes and match the determinant ratio at $P_0$.

Finally, the clean-extension argument lifts as well. In its product
presentation $(g,y)$, the ratio between the target and source determinant
lines is the multiplicative line
$\mathcal R(g)=\det(g\Lambda:\Lambda)$. The splitting on the
conjugated congruence group descends this line through the quotient used
in \cite[Appendix~A.2]{GaitsgoryWhit}, and the unipotent splitting gives
its left $N$-action. The ratio line can remain nontrivial on the boundary;
only its equivariance is needed. On every finite-dimensional boundary
piece the proof supplies a unipotent
stabilizer on which the Whittaker character is nontrivial
\cite[Appendix~A.2--A.3]{GaitsgoryWhit}. This stabilizer fixes the source coordinate and acts on the ratio-line
fiber by an algebraic character to $\mathbb G_m$, hence trivially. It
therefore fixes the target frame as well, while retaining its nontrivial
Whittaker character. The same boundary-vanishing argument
applies to arbitrary objects on the frame torsor. Together with the
lifted averaging correspondences, this transports the units, counits,
and clean-extension isomorphisms of the unit-restriction equivalence.

Every lifted map commutes with multiplication of the determinant frame by
$\mathbb G_m$. Taking strong equivariance with the Kummer character of
exponent $s$ therefore descends the lifted equivalence to D-modules
twisted by $\mathcal L^s$. An equivariant equivalence induces an
equivalence on these descent categories, including their full
cosimplicial descent data. This yields the stated determinant-level
comparison, and preserves its base and Ran compatibilities.
\end{proof}

We will need two point inputs. At an irrational scalar level, the point
Whittaker Grassmannian category is a sum of copies of
$\operatorname{Vect}$ indexed by the dominant coweights of the actual
group. Here is the precise consequence of Campbell--Dhillon--Raskin that
we use. The irrational case of their integral-Weyl-group calculation
\cite[Theorem~3.5.6, Equation~(3.28)]{CDR} leaves one dominant label in
each contributing spherical highest-weight block. Its finite-length
one-simple highest-weight heart is therefore finite-dimensional vector spaces.
The derived-heart comparison of \cite[Section~2.9.4]{CDR}, together with
the compact generators used in the proofs of their Corollary~3.5.9 and
Theorem~3.7.3, identifies the presentable block with
$\operatorname{Ind}(D^b(\operatorname{Vect}^{\mathrm{fd}}))=
\operatorname{Vect}$. Thus higher self-extensions are excluded by the
derived comparison, not just by the absence of linkage. Their spherical
fundamental local equivalence \cite[Remark~3.7.5]{CDR} transfers this
statement to Whittaker D-modules. Their actual-form reduction
\cite[Section~4.2, Proposition~4.2.2, and Section~4.3]{CDR} retains the
actual weight and coweight lattices and treats products. An irrational
complex level of either sign is in their negative good-level regime.
For groups with simply connected derived group, the corresponding
irrational clean-standard statement was already proved in
\cite[Proposition~8.4]{GaitsgoryTwisted}.

At exponent zero we use the clean-standard theorem
\cite[Theorem~3]{FrenkelGaitsgoryVilonen}; its Introduction explicitly
includes characteristic-zero D-modules. Cleanliness and the mapping
complex calculation in Proposition~\ref{whittaker:standards} give the
corresponding DG semisimplicity. This is the categorical form of geometric
Casselman--Shalika needed here. When passing between invariant and
coinvariant conventions at a point, we use the Whittaker
invariant/coinvariant theorem \cite[Theorem~2.1.1]{RaskinW}. These are point
statements; the family conclusions used below will be deduced separately.

\subsection{Defect strata and a continuous Whittaker projection}

In a fixed pole bound, record actual poles and defects by
\[
 D=\sum_y\lambda_y y,\qquad T_D=T_0(-D),
\]
with poles positive. Orders of Pl\"ucker maps pair additively with dominant
weights. Indeed their leading terms have nonzero Cartan product, so their
valuations add. Using every highest weight of $H$ shows that each
$\lambda_y$ is a coweight of $H$. Away from the markings it belongs to the
negative rational coroot cone. At a marking it is bounded above by the
specified pole bound, in the order defined by that cone.

Stratify both the markings by their coincidences and the Pl\"ucker maps by
this divisor. Write $S_D$ for the resulting divisor-parameter scheme and
\[
 f_D:\mathcal Y_D\hookrightarrow\mathcal Y_m,\qquad
 q_D:\mathcal Y_D\longrightarrow S_D
\]
for the inclusion and the parameter map. Dividing the maps by their common
zeros gives a genuine Borel reduction with Cartan bundle $T_D$.
Conversely such a reduction recovers the point of the stratum. Thus
$q_D$ parametrizes $N^{T_D}$-bundles. It is a smooth fibration given by a
tower of vector stacks: filtering $N$ by root height reduces each stage
to the stack associated with the two-term complex of global sections of
a root line bundle. The determinant twist restricts to a line from $S_D$,
by the same filtration.

This description is valid in families. After clearing bounded poles,
flattening by the lengths of zero divisors makes them relative Cartier
divisors. One then refines by coincidences, and, when needed, passes to
\'etale charts enumerating the support. D-modules allow us to ignore
nilpotents in this stratification. Opens bounding the total zero lengths
exhaust $\mathcal Y_m$; over quasicompact parameter pieces they are
quasicompact and have only finitely many of these strata. Their
stabilizers are unipotent. In fact an automorphism preserving all
Pl\"ucker maps restricts generically to the unipotent stabilizer of a
reduction with fixed Cartan datum. Alternatively, on a bounded open a
sufficiently deep frame at a nondefective point kills the stabilizer.

A stratum is \emph{relevant} if every $\lambda_y$ is dominant. Dominance
and the negative-cone condition force all off-marking defects to vanish.
There are only finitely many relevant labels in each pole bound.
On a relevant stratum the simple-root lines are
\[
 \alpha(T_D)=\Omega_X\bigl(-\langle\alpha,D\rangle\bigr)
 \longrightarrow\Omega_X.
\]
The simple-root abelianization and the trace on $H^1(X,\Omega_X)$ define a
function $\psi_D:\mathcal Y_D\to\mathbb A^1$. Choose its sign so that
modification at a nondefective point changes $\psi_D$ by the residue
character of the modifying loop. In Grassmannian coordinates its phase
along $nt^\lambda$ is therefore $\chi(n)$.

\begin{proposition}[Projection and stratum tests]
\label{whittaker:strata-properties}
In every pole bound, inclusion of the strongly Whittaker category has a
continuous right adjoint $\operatorname{Av}_*$. Its image consists exactly
of the objects $F$ satisfying
\begin{enumerate}
\item $f_D^!F=0$ on every irrelevant stratum;
\item on a relevant stratum,
$f_D^!F=\operatorname{Exp}(\psi_D)\otimes^!q_D^!E_D$ for an object $E_D$
in the D-module category of $S_D$ with its induced twist.
\end{enumerate}
The analogous statement uses $-\psi_D$ for the opposite character.
Restriction by $!$ computes the projection on each stratum. These
projections commute with $!$-base change in the parameter schemes and
with tensoring a spectator category. They are compatible with
$!$-restriction between pole bounds. Closed direct images, both between
pole bounds and along closed parameter embeddings, preserve the
Whittaker subcategories.

If a bounded-zero open contains every relevant stratum of a pole bound,
restriction to that open and star extension identify its Whittaker
category with that of the entire bound.
\end{proposition}

\begin{proof}
We first construct the projection on a quasicompact bounded-zero open on
which an auxiliary point $z$ is unmarked and nondefective. Use root-height
coordinates for $N_0$ at $z$. Let $U_n$ be the unipotent loop subgroup
allowing pole order at most $n\operatorname{ht}(\alpha)$ in the
positive-root coordinate $\alpha$; $U_0$ is the arc subgroup. Root heights
add under multiplication, so these are subgroups. Each $U_n$ is
pro-unipotent and the homogeneous space $U_n/U_0$ is finite dimensional
and de Rham acyclic; $U_0$ need not be normal in $U_n$.

Let $\mathcal Z^{\mathrm{fr}}$ be the formally framed stack on this
open. It carries a $U_n$-action, and the unframed open is
$\mathcal Z^{\mathrm{fr}}/U_0$. Its averaging arrows form the \v Cech
groupoid of
\[
 \mathcal Z^{\mathrm{fr}}/U_0
   \longrightarrow\mathcal Z^{\mathrm{fr}}/U_n.
\]
Equivalently, these arrows are isomorphisms away from $z$ and
modifications by $U_n$ at $z$. The character descends to the groupoid
because $\chi|_{U_0}=0$; normality of $U_0$ is unnecessary. Write
$s_n,t_n$ for source and target. They are smooth, with
fiber $U_n/U_0$ of dimension $d_n$, and preserve the pole bounds and defect
strata. Indeed the Pl\"ucker maps agree away from $z$, and $N$ fixes
every highest-weight vector at $z$; the modified reduction is therefore
still nondefective there. Thus the actual defect divisor and the
bounded-zero open are preserved. If finite-type charts are desired, divide frames by an arc subgroup
sufficiently deep to be normal in $U_n$. The arrow splitting of the
determinant line defines the averaging functor
\begin{equation}\label{eq:W1}
 \operatorname{Av}_{n,*}(F)=
 t_{n,*}\bigl(\operatorname{Exp}(\chi)\otimes^!s_n^!F\bigr)[-2d_n].
\end{equation}
Smooth pullback along the unipotent fibers is fully faithful because
their de Rham cohomology is the ground field. Base change and additivity of
$\chi$ identify \eqref{eq:W1} with a strongly equivariant object for this
finite groupoid. Smooth star-pullback/star-image adjunction then identifies
it with the right adjoint to inclusion of that equivariant subcategory.
In particular its counit and the transition maps for larger $n$ are
canonical. The maps in \eqref{eq:W1} are finite-type and safe after taking
these quotients, so the functor is continuous on all D-modules.

We next compute the restrictions of \eqref{eq:W1}. On a defect stratum its
arrow stack maps to $\mathcal Y_D\times_{S_D}\mathcal Y_D$. For large $n$
this map is, smoothly locally, a tower of affine-space bundles. To see
this, lift isomorphisms between the two unipotent extensions in order of
root height. At the $\alpha$-stage, possible lifts form a torsor under
global sections of the corresponding root line with poles
$n\operatorname{ht}(\alpha)z$. The obstruction lies in its $H^1$ and
vanishes for large $n$. Apply relative Serre vanishing to the bounded
family of root lines on $X\times S_D$. Once $R^1p_*$ vanishes,
Riemann--Roch and cohomology-and-base-change make the spaces of sections
locally free of constant rank, compatible with arbitrary parameter base
change. Hence their torsors are smoothly locally affine-space bundles
even where the unmodified $H^0$ and $H^1$ jump. Higher-root lifts are
controlled by the same graded root lines. The height convention makes these successive extensions
compatible with multiplication.

On a relevant stratum the character on the arrow stack is
$t_n^*\psi_D-s_n^*\psi_D$, by the residue description of the trace.
Base change and de Rham integration along the preceding affine-space
towers therefore give
\begin{equation}\label{whittaker:stratum-projector}
 F\longmapsto\operatorname{Exp}(\psi_D)\otimes^!
 q_D^*q_{D,*}\bigl(\operatorname{Exp}(-\psi_D)\otimes^!F\bigr).
\end{equation}
For the shift, put $e=\dim(q_D)$ and let $h=d_n-e$ be the relative
dimension of the affine-space tower from the arrow stack to
$\mathcal Y_D\times_{S_D}\mathcal Y_D$. Integration of its $!$-pullback
contributes $[2h]$, so
\[
 [2h-2d_n]=[-2e],\qquad q_D^![-2e]=q_D^*.
\]
The two pullbacks have the same essential image, differing by this fixed
shift on the stratum. Pullback is fully faithful for the vector-stack tower
$q_D$, so \eqref{whittaker:stratum-projector} is its exponential-conjugate
projection.

On an irrelevant stratum there is a simple root $\alpha$ for which
$\alpha(T_D)$ admits a pole at some point other than $z$. A section with
sufficient pole allowance at $z$ can prescribe a nonzero residue at that
point and zero residues at the other allowed poles. The required
surjectivity is again the vanishing of $H^1$ after increasing $n$.
The residue theorem makes its residue at $z$ nonzero. Thus, after
integrating the higher-root coordinates, the averaging contains the de
Rham cohomology of a nonconstant linear exponential on an affine space;
this cohomology is zero. The calculation works over the parameter scheme,
\'etale locally if a support enumeration is needed.

There are finitely many strata in our open. Choose $n$ large enough for
all of them. On every stratum the adjunction transitions from larger
averaging functors to \eqref{eq:W1} are now isomorphisms: on relevant strata
all are the projection \eqref{whittaker:stratum-projector}, with its
canonical counit, and on irrelevant strata all are zero. The finite
stratification is conservative for $!$-restriction, so the transitions are
isomorphisms on the open. Equivalently, the counits testing equivariance
for larger $U_n$ become isomorphisms there. The stable functor is therefore
the right adjoint to the full Whittaker inclusion, with exactly the two
stated stratum tests.

All calculations were relative and used finite-type correspondences, so
they commute with $!$-base change and spectators before stabilization.
The same uniform finite choice proves that they still do so afterwards.
On overlaps of auxiliary-point opens the tests describe the same full
subcategory. The right adjoints and their counits consequently glue,
uniquely compatibly with the inclusions. Several or varying auxiliary
points give the same condition: repeat the calculation over their
parameter scheme, replacing $nz$ by the corresponding auxiliary divisor
and adding their characters. This also verifies the smooth-local version
of strong equivariance.

The bounded-zero opens exhaust the pole bound, and their projections glue
by these compatibilities. Closed images preserve the two tests by
$!$-base change, both for an increase of pole bound and for a closed
parameter embedding. Finally, outside an open containing all relevant
strata the $!$-restriction of a Whittaker object is zero. The localization
triangle identifies it with star extension from that open. Conversely
star extension preserves the tests and hence recovers a Whittaker object.
This proves the last assertion and completes the proof.
\end{proof}

\subsection{Duality and passage to the Ran space}

The projections just constructed allow duality to be established on the
full Whittaker categories. The opposite character is essential: it
cancels the exponential in the pairing.

\begin{proposition}[Whittaker duality]\label{whittaker:duality}
The categories $\mathcal W_s^+(S)$ and $\mathcal W_{-s}^-(S)$ are dual,
with the normalized pairing
\begin{equation}\label{eq:W2}
 \langle W,W'\rangle=
 \Gamma_{\mathrm{dR}}(W\otimes^!W')[-2\delta_N].
\end{equation}
In a pole bound this formula uses the bounded-zero star-extension
convention of Proposition~\ref{whittaker:strata-properties}. It then
extends continuously to the full pole model. Equivalently, one may place
one argument in a pole bound and $!$-restrict the other to that bound.
The analogous pairing relative to $S$ is defined by relative star image.
Closed images transpose to $!$-restrictions.

These statements hold for the Ran categories, formed as limits under
$!$-restriction along the diagonals, and are compatible with tensoring a
spectator category.
\end{proposition}

\begin{proof}
First work in a bounded-zero quasicompact open of a pole bound. Its
unipotent inertia makes ordinary de Rham pushforward safe. The twisted
QCA duality pairing is therefore ordinary de Rham global sections after
the twists cancel.

Here is a direct way to pass the smooth-stack duality of
Proposition~\ref{foundations:categories} to these pole stacks. On a
quasicompact open of $\operatorname{Bun}_H$, allow sufficiently many extra
poles that the relevant $H^1$ groups vanish. The spaces of Pl\"ucker maps
are then vector bundles, and their relations are closed conditions in a
smooth quotient stack. An open condition can be realized by restricting
the ambient stack to an open, so the bounded piece admits a closed
embedding into such a smooth ambient stack, with the twisting line
extended from it. For the parameters, work on affine opens embedded as closed subschemes
of smooth ambient schemes, using the graph of the marked-point map and
local trivializations of the twisting line. For a closed embedding $i$, Kashiwara's equivalence realizes the
supported category as the image of the continuous projector $i_*i^!$.
On a product the two external projectors commute and cut out precisely
the product support. Thus K\"unneth and the diagonal duality identities
pass to these closed supports. After cancellation of the twists,
renormalized pushforward in those identities agrees with ordinary
pushforward by safety. This supplies the stated pairing, also with
spectators.

Let $P_+$ and $P_-$ denote the projections for the two opposite characters.
For a Whittaker object $W$ and an arbitrary oppositely twisted object $F$,
the counit $P_-F\to F$ induces an isomorphism
\begin{equation}\label{whittaker:orthogonal-projection}
 \Gamma_{\mathrm{dR}}(W\otimes^!P_-F)
 \xrightarrow{\ \sim\ }
 \Gamma_{\mathrm{dR}}(W\otimes^!F).
\end{equation}
Indeed, on an auxiliary-point open insert the finite formula
\eqref{eq:W1} for $P_-$ with $n$ sufficiently large. The $!$-projection
formula moves $W$ to the arrow stack. Its positive-character equivariance
cancels the negative exponential in the averaging kernel. Integrating
along the source then cancels the dimension normalization of
\eqref{eq:W1} and gives the right side of
\eqref{whittaker:orthogonal-projection}. The identification is the one
induced by the counit. Open descent gives it on the whole quasicompact
open. The same argument interchanges the two characters.

Consequently the Whittaker subcategory on either side annihilates the
kernel of the projection on the other. Transposing the inclusion and its
continuous right adjoint under ambient duality gives adjoints in the
opposite order. The transposed projector is the identity on the opposite
Whittaker image and zero on its projection kernel, by
\eqref{whittaker:orthogonal-projection} and its counterpart. It is
therefore exactly the opposite projector. Splitting these projectors
identifies the two Whittaker categories as duals. The fixed shift
$[-2\delta_N]$ is an invertible normalization of this duality.

An open containing all relevant strata computes the Whittaker category of
the entire pole bound. The tensor product in \eqref{eq:W2} is likewise
star-extended from that open: its $!$-restriction to the complement is
zero. This proves the formula and duality on the bound.

For closed increases of pole bound, the continuous closed image is left
adjoint to the continuous $!$-restriction. The latter also preserves
limits. The limit of these right adjoints is the colimit of their left
adjoints in presentable categories. Formula~\eqref{eq:W2}, together with
the $!$-projection formula, transposes the closed image to the
corresponding restriction. Passing to this limit/colimit therefore
preserves the duality.

The Ran presentation has the same form. By
Proposition~\ref{whittaker:strata-properties}, its diagonal restrictions
and closed images are computed in the collision parameter schemes, after
a cofinal increase of pole bounds. Their adjunctions and transpose
identities are the ones just proved. Thus
\[
 \mathcal W_s^+(\operatorname{Ran})
   =\lim_{I,\,\Delta^!}\mathcal W_s^+(X^I)
   =\operatorname*{colim}_{I,\,\Delta_*}\mathcal W_s^+(X^I),
\]
with the corresponding statement for the opposite character. To see the
duality directly, compute maps out of the colimit by the diagram and use
the termwise dualities; the resulting dual is exactly the displayed
limit on the opposite side. Tensoring by a spectator category commutes
with this presentation because it commutes with the colimit, and the
continuous adjunctions tensor with it. The same argument gives the
relative pairing over $S$.
\end{proof}

\subsection{Standard objects and the t-structure}

We now specialize to an irrational exponent $s$ or to $s=0$. The point
semisimplicity recalled above makes the relevant strata into actual
summands. We must establish this over parameter schemes, since a
pointwise statement alone would not control arbitrary D-modules in a
family.

Let $S$ be a fixed coincidence piece, so that the distinct marked points
are indexed by a fixed finite set. A relevant label is a tuple of
dominant coweights at these points. For such a label $D$, the scheme
$S_D$ is identified with $S$; write $\mathcal L_D$ for the line on $S$
whose pullback is the determinant line on $\mathcal Y_D$. Set
\[
 a_D=\sum_y\langle 2\rho_{\mathrm{rt}},\lambda_y\rangle,
 \qquad 2\rho_{\mathrm{rt}}=\sum_{\alpha>0}\alpha.
\]
Riemann--Roch, applied to the root-height filtration, gives
\begin{equation}\label{whittaker:stratum-dimension}
 \dim(q_D)=-\chi(X,\mathfrak n^{T_D})=\delta_N+a_D.
\end{equation}

\begin{proposition}[Standard summands in families]
\label{whittaker:standards}
Suppose $s$ is irrational or zero. Over a fixed coincidence piece $S$,
the Whittaker category is the direct sum, in presentable DG categories,
of the categories $D(\mathcal L_D^s)$ indexed by relevant labels $D$.
The inclusion of the $D$-summand is star extension from its stratum;
with the normalization below it is
\begin{equation}\label{eq:W3}
 E\longmapsto
 f_{D,*}\bigl(\operatorname{Exp}(\psi_D)\otimes^!q_D^!E[-a_D]\bigr).
\end{equation}
The construction takes place in any pole bound containing the label and
is then carried to larger bounds by closed image. For holonomic
arguments in this formula, the canonical shriek-to-star map is an
isomorphism. Both assertions hold for the opposite twist and character.
\end{proposition}

\begin{proof}
Begin over a point. The star standard and its shriek counterpart are
Whittaker objects. For the former this follows from the stratum tests;
for the latter use the smooth modification arrows, equivariance, and
partial shriek on the holonomic exponential. Their endomorphism
complexes are the ground field: locally closed image is fully faithful,
and pullback along the unipotent stack $q_D$ is fully faithful.

For irrational $s$, the canonical map from shriek to star is nonzero,
since its restriction to the stratum is the identity. In the semisimple
point category supplied by the block theorem and
Proposition~\ref{whittaker:model-comparison}, an object with scalar
endomorphism complex is a shift of a single simple. The nonzero map
between these two objects is therefore an isomorphism. For $s=0$, the
same cleanliness is precisely the imported clean-standard theorem.
In either case, adjunction and cleanliness make the mapping complexes
between distinct standard objects zero: restriction of one clean
standard to the other stratum is zero. Moreover, partial shriek
adjunction against an arbitrary Whittaker object $W$ identifies maps
from the $D$-standard with the coefficient complex extracted from
$f_D^!W$ by exponential cancellation and the fully faithful $q_D^!$,
up to the displayed fixed shift. This extraction is continuous by
Proposition~\ref{whittaker:strata-properties}. Each standard is therefore
compact in the Whittaker category. The standards generate by the stratum
tests and finite stratification in each bound. These orthogonal compact
generators with scalar endomorphisms give the asserted presentable DG
summand description, including at exponent zero.

Over a fixed coincidence piece choose, \'etale locally, curve coordinates
near the distinct marked sections and compatible spin frames. If $z$ is
such a coordinate, $t=z-z(x_i)$ is an actual relative parameter; formal
\'etaleness identifies the completed graph neighborhood with
$\mathcal O_S[[t]]$. Thus no section of an arbitrary formal-coordinate
torsor is required. The spin frames give the corresponding Cartan
trivializations. The local Grassmannians, residue
characters, loop actions, and determinant twists are then products of
the point models with the parameter scheme, with the possible residual
line on the parameters retained. The point decomposition tensors with
$D(S)$. Taking invariants of this product commutes with that tensor
product because $D(S)$ is dualizable. This proves the family decomposition
in these coordinates, with all D-modules as coefficients.

To identify and descend its summands, take $!$-restriction along the
sections $t^\lambda$ in the Grassmannian model. Under the model comparison
this is evaluation of the corresponding exponential-conjugate $q_D^!$
object on the global label stratum, carrying exactly the twist
$\mathcal L_D^s$. The standard defined by $f_{D,*}$ has these restrictions
by base change. Hence the summands and their inclusions are intrinsic,
and descend under changes of coordinates and twisting-line
trivializations. This gives the asserted formula before its harmless
normalizing shift.

For a holonomic parameter object, shriek extension of the stratum object
is Whittaker by the same smooth-arrow calculation. Partial adjunction
computes its maps to Whittaker objects from their $!$-restrictions to the
label stratum. The summand description shows that star extension
represents the same functor. The canonical shriek-to-star map is the
resulting isomorphism. This proves the family cleanliness assertion.
\end{proof}

Use the usual stack perverse normalization on D-modules: $!$-pullback to
a smooth chart is shifted back by its relative dimension. On the pole
model translate the resulting t-structure so that its heart is the stack
heart shifted by $[\delta_N]$. This fixes the Whittaker t-structure used
from now on.

\begin{proposition}[Normalization and divisor amplitude]
\label{whittaker:t-structure}
This t-structure restricts to each Whittaker pole bound and extends to the
full pole model with t-exact closed transition images. With this
convention \eqref{eq:W3} is t-exact. On an affine chart where a smooth
divisor diagonal has inclusion $i$ and affine open complement $j$,
\[
 i_*,\ j^!,\ j_*\quad\text{are t-exact},
 \qquad i^!\quad\text{has t-amplitude }[0,1].
\]
The resulting t-structures on the parameter categories used here are
separated.
\end{proposition}

\begin{proof}
The strong character condition is preserved by truncation: its
isomorphisms are tested by the t-exact smooth pullbacks normalized by
relative dimension. Closed transition images are t-exact in the ambient
D-module categories and hence on their Whittaker subcategories.

An object in the full pole model is the filtered colimit of the closed
images of its $!$-restrictions to pole bounds. Truncate these objects in
the bounds. The transition maps between their truncations are compatible
because the closed images are t-exact. This constructs the two
truncation terms in the colimit category. Orthogonality follows by
adjunction: restriction to any fixed bound is left t-exact and commutes
with filtered colimits. The resulting t-structure is compatible with
filtered colimits.

On a fixed coincidence piece, Proposition~\ref{whittaker:standards}
decomposes a pole bound into summands, and its t-structure decomposes
with them because its aisles are closed under direct summands. For one
label, work in an open where its stratum is closed. Open pullback and
closed image are t-exact. That open pullback is conservative on the
summand, so the t-structure is detected there. The smooth normalization
and \eqref{whittaker:stratum-dimension} now give the precise shift:
$q_D^![-(\delta_N+a_D)]$ takes a parameter heart object to the stack
heart, and translating the latter by $[\delta_N]$ gives exactly
$q_D^![-a_D]$. Exponential tensoring preserves this normalization. This
proves the t-exactness of \eqref{eq:W3}; closed transition images give it
in the entire pole model.

For a smooth divisor and its principal affine complement, the asserted
exactness is the ordinary D-module exactness for a closed inclusion and
an affine open inclusion, applied on the ambient charts. The localization
triangle then places $i^!$ in amplitude $[0,1]$. Base change with closed
images and filtered colimits passes these assertions to the pole
presentation. Finally, $!$-restrictions to the finitely many parameter
coincidence strata have bounded t-amplitude and jointly detect objects.
On each such stratum the standard decomposition reduces separation to
that of D-modules on a scheme. This proves separation in the parameter
categories at issue without an assertion about arbitrary ind-scheme
t-structures.
\end{proof}

\subsection{The coefficient functor}

We can now define the continuous Whittaker coefficient functor with all
its required parameter operations:
\begin{equation}\label{whittaker:coefficient-functor}
 R_H:D(L_H^s)\longrightarrow\mathcal W_s^+(\operatorname{Ran}).
\end{equation}
At marked parameters it is $p_{\mathcal Y}^!$ followed by
$\operatorname{Av}_*$ of Proposition~\ref{whittaker:strata-properties}.
The base-change and closed-image compatibilities of that proposition
assemble the marked-point functors into \eqref{whittaker:coefficient-functor}.
We use the analogous functor for $-s,-\chi$, and the same definition at
$s=0$. Proposition~\ref{whittaker:duality} supplies the duality for these
actual presentable Whittaker categories. Thus the adjoints and transposes
used below are defined on all D-modules in the base, while shriek
operations on holonomic standards serve only to establish their stated
cleanliness.
